L'attivazione del flusso rende visibili le variazioni dei parametri meccanici, ma la loro osservabilità dipende dalla distribuzione di $c$ e dalla risposta alla sorgente $-\nabla\cdot(c\,\delta\mathbf q)$ della~\eqref{eq:mech-tracer-variation}. Nelle prove sul quadrato il rango pieno compare quando la quarta singolarità è ancora piccola. Il Péclet globale caratterizza la famiglia di carichi studiata. La collocazione dei sensori e la direzione della perturbazione di flusso determinano quanta parte del segnale arrivi alle misure.

Lo stesso meccanismo chiarisce il risultato sui sensori. Il condizionamento globale confronta le direzioni estreme del Jacobiano, mentre $s_{\beta\mid-\beta}$ misura quanta variazione del tracciante associata a $\beta$ sopravvive alla compensazione. I profili FEM confermano che possono coesistere rango pieno e risposte quasi indistinguibili a ordini meccanici diversi. Il Monte Carlo quantifica questa fragilità nel disegno bilanciato a quattro incognite. Le repliche linearizzate a sei incognite illustrano l'effetto dei vincoli in un protocollo distinto.

Con etichette per cella, l'inversione assorbe parte della discrepanza fra discretizzazioni modificando i coefficienti. Con il raster comune, la discrepanza diretta a 3--2~mm resta comparabile in norma al contrasto fra M2 e M4 a parametri fissati. Per stabilire il bias su $\beta$ serve la componente dello scarto lungo la sensibilità residua $r_\beta$ in~\eqref{eq:beta-discretization-bias} e una verifica con inversioni non lineari a raster comune.

Il dominio BraTS è una singola sezione, con forzanti e concentrazioni sintetiche. Le scale del calcolo servono a definire il confronto numerico. La diffusività nominale, per esempio, è maggiore dei valori preclinici riportati per il Gd-DTPA~\citep{he2013gdtracer,wang2019gdtracer}. L'estensione a misure longitudinali e coefficienti calibrati richiede una validazione distinta.
